% Chapter Template

\chapter{INTRODUCTION} % Main chapter title

\label{Chapter1} % Change X to a consecutive number; for referencing this chapter elsewhere, use \ref{ChapterX}

\lhead{Chapter 1. \emph{INTRODUCTION}} % Change X to a consecutive number; this is for the header on each page - perhaps a shortened title

%----------------------------------------------------------------------------------------
%	SECTION 1
%----------------------------------------------------------------------------------------

\section{Introduction}
Natural and man-made disasters such as floods, cyclones, earthquakes, landslides, forest fires, and industrial accidents frequently damage the very infrastructure that people rely on to communicate. Cellular towers lose power or are physically destroyed, fibre and microwave backhaul links are severed, and even the networks that survive are often overwhelmed by the sudden surge of calls and messages. The first hours after a disaster are the most critical for search and rescue, yet they are exactly the hours in which communication is most likely to be unavailable. Existing alternatives such as satellite phones and HF/VHF radios have their own limitations: satellite phones are too expensive to place in the hands of every responder or village, while conventional handheld radios have limited range, offer no location tracking or message logging, and depend on trained operators. This gap has led to growing interest in low-power wide-area network (LPWAN) technologies for emergency use, and among these, LoRa (Long Range) stands out as a particularly promising option. LoRa is a radio modulation technique based on chirp spread spectrum that operates in licence-free sub-GHz ISM bands, such as 865-867 MHz in India. It can cover several kilometres in urban surroundings and well over ten kilometres in open, line-of-sight terrain, while consuming so little power that a node can run for weeks or months on a small battery or a modest solar panel. Its high receiver sensitivity allows signals to be decoded even when they are weaker than the background noise, and its adjustable spreading factors let designers trade data rate against range. The price of these strengths is low throughput, which makes LoRa unsuitable for voice or video, but ideal for the short text messages, SOS alerts, GPS coordinates, and sensor readings that matter most in a crisis. This project proposes a LoRa-based emergency communication network that operates independently of cellular and internet infrastructure. It is built from portable battery-powered end nodes carried by responders and residents, solar-assisted relay nodes placed on high ground or buildings to extend coverage through multi-hop paths, gateways that collect field traffic and forward it onward when any backhaul is available, and a command-center dashboard that displays node locations, messages, and alerts on a live map. Together these elements form a self-healing, scalable, and low-cost network that can be deployed quickly where conventional communication has failed.
\section{Objectives}
The general objective of the project is to design, implement, and evaluate a LoRa-based communication network that delivers reliable short-message and location services under emergency conditions where conventional communication infrastructure is unavailable. To achieve this, the project sets out several specific technical objectives. First, it aims to build prototype end nodes, relay nodes, and a gateway using LoRa transceivers, microcontrollers such as the ESP32 or STM32, GPS receivers, and suitable antennas, and to develop firmware that handles message composition, SOS triggering, GPS acquisition, acknowledgment, retransmission, and power management. Second, it aims to design a lightweight network protocol covering addressing, packet structure, routing by controlled flooding or a mesh and store-and-forward scheme, and collision avoidance, so that messages can travel across several relays and coverage can extend far beyond the range of a single link. Third, it seeks to make the network reliable through acknowledgments, retries, duplicate suppression, and message prioritization, giving SOS messages the highest priority, while keeping energy use low through sleep modes, adaptive transmission power and spreading factor, and solar charging. Fourth, it intends to secure the network by encrypting payloads with AES and authenticating nodes to guard against spoofing, replay, and tampering. On the functional side, the project aims to let any node send a one-press distress alert carrying GPS coordinates, node identity, timestamp, and battery status; to support two-way messaging between nodes and the command center, including predefined quick messages that are easy to use under stress; and to develop a dashboard that presents a live map, a message feed, alerts, and node health indicators. It further aims to ensure that the field network and gateway work fully offline, storing data locally and synchronizing later, and to forward critical alerts through SMS, email, or internet services whenever any uplink becomes available. Finally, the evaluation objectives are to measure range, packet delivery ratio, end-to-end latency, throughput, and battery life across different terrains and conditions, to determine how many nodes the network can support before congestion degrades performance, to document the per-node and per-deployment cost to demonstrate affordability, and to test the interface with non-technical users to confirm that it is simple enough to use in a real emergency.

\subsection{Purpose}
The purpose of this project is to create a dependable, independent, and affordable communication layer for emergencies, one that preserves the ability to call for help, coordinate a response, and share situational information at the moment when every other channel has failed. For affected individuals and communities, the system offers a simple way to raise an SOS with accurate location even where there is no mobile signal, to report their condition so that relief can be prioritized, and to reduce the isolation and panic that come with sudden loss of contact. For rescue and response teams, it maintains a link between field personnel and the control room, allows the movements of teams to be tracked for their own safety, and supports faster, better-informed decisions about where to send people and supplies. For disaster-management authorities and local governments, it provides a low-cost, locally owned system that can be pre-positioned in vulnerable regions, that records messages and movements for later review, and that complements rather than replaces existing disaster-management systems. The project also serves an educational and research purpose, since it demonstrates a practical application of embedded systems, IoT, and wireless networking for social benefit and provides an open, extensible platform for further work on mesh routing, energy harvesting, and resilient network design.
\subsection{Scope}
The scope of the project covers the full path from design to field validation of a prototype network. On the design side, it includes the overall architecture and topology, the link budget analysis and frequency plan, the selection of LoRa parameters such as spreading factor, bandwidth, and coding rate, and guidelines for antenna selection and placement. On the hardware side, it includes selecting and integrating LoRa modules, microcontrollers, GPS modules, displays, buttons, and buzzers; designing the power subsystem with rechargeable batteries, charge controllers, and solar panels for relays; building weather-resistant enclosures; and preparing the gateway, typically a single-board computer with a LoRa radio or concentrator. On the software side, it includes the embedded firmware for end and relay nodes, the packet formats and routing logic, the gateway services that collect, store, and forward data, the command-center dashboard with mapping, alerts, message history, and basic analytics, and optionally a smartphone application that pairs with a node over Bluetooth to make text entry easier. Security measures such as payload encryption and key management for the prototype fall within scope, as does a complete testing program consisting of laboratory tests of radio performance and power consumption, network simulation to study scalability and routing behaviour, outdoor trials in open, semi-urban, and obstructed terrain, and mock disaster drills that exercise the workflow from distress call to response. The deliverables include design documents, schematics, source code, a deployment guide, a user manual, and a final test report. Several matters fall outside the scope of this phase of work: voice, video, and image transmission or any application requiring high bandwidth; mass manufacturing, regulatory type approval, and commercial certification; integration with national or state disaster-management platforms, which is left as future work; the design of satellite or cellular networks, which are considered only as optional backhaul; custom silicon or radio design; and long-term operation, maintenance, and training programs beyond the prototype stage. The work rests on certain assumptions, namely that the system operates in licence-free ISM bands within the legal limits of the country of deployment, that messages remain short at tens of bytes, that nodes generally have GPS visibility although indoor or shadowed nodes may report their last known position, that some relays can be placed at elevated sites, and that solar or battery power is available at relays and the gateway. It also acknowledges inherent limitations: LoRa’s low data rate restricts the type of content that can be carried, actual range varies with terrain, vegetation, buildings, weather, and antenna height, capacity is limited because many simultaneous transmissions increase collisions, regulatory rules cap transmit power and sometimes duty cycle, battery life depends on reporting intervals and the power-hungry GPS, and the budget and timeline of a prototype restrict the number of nodes that can be built and tested.

\subsection{Applicability}
The network is applicable wherever conventional communication is unavailable, damaged, or overloaded. In disaster and emergency management, it can coordinate boats, shelters, and relief camps during floods and cyclones when towers are down, allow rescuers and trapped survivors to signal their positions after earthquakes and building collapses, support communication in the mountainous terrain of landslides and avalanches where coverage is already poor, keep firefighting crews in contact and tracked during forest fires, and help coordinate evacuation during industrial accidents or chemical leaks when normal networks are disrupted. It is equally valuable in remote and underserved regions such as rural villages, hill communities, tribal and forest areas, border regions, and islands, where it can also form the basis of community-owned early-warning and alert systems. In coastal and maritime settings it can carry distress signals from fishermen and small boats operating beyond the reach of towers, deliver weather alerts to registered nodes, and help coastal rescue teams coordinate. It also suits safety and tracking applications for trekkers, mountaineers, and pilgrims, for wildlife reserves and forest-department patrols including anti-poaching teams, and for mining, tunnelling, and construction sites where signals are blocked and worker safety is a concern. Large public gatherings such as festivals, rallies, and sporting events, where mobile networks often become congested, can use it as a backup channel, and hospitals, utilities, and police can rely on it during outages, as can relief camps and temporary field hospitals. Defence and civil-defence organizations may use it as a quick-deploy field network for relief operations and training drills, and educational institutions can use it as a hands-on platform for teaching IoT and wireless networking. The principal beneficiaries are affected citizens, who gain the ability to call for help and report their status without a phone network; first responders such as NDRF and SDRF teams, fire services, police, and medical staff, who gain coordination, tracking, and situational awareness; disaster-management authorities, who gain a low-cost, deployable backup; NGOs and volunteer groups, who gain affordable tools for relief coordination; local governments and communities, who gain locally owned alert infrastructure; and students and researchers, who gain an open platform for further work. Compared with existing solutions, the system is far cheaper per node than satellite phones or professional radio systems, runs for long periods on batteries and solar power, needs no operator network or internet, scales simply by adding nodes and relays, provides logged and location-aware data rather than voice alone, and can be extended to carry sensor data such as water level, temperature, smoke, or seismic activity for early warning.
